Die Suche nach einem geeigneten Platz für das Praxisprojekt begann im Sommer 2025 mit dem Ziel, ein Unternehmen zu finden, das eine intensive Auseinandersetzung mit der reinen 
Programmierung ermöglicht. Der entscheidende Kontakt zur snoopmedia GmbH kam über ein persönliches Netzwerk zustande: Ein Kommilitone, mit dem ich bereits im Rahmen des Moduls 
„Software Engineering 2“ zusammengearbeitet hatte, war zu diesem Zeitpunkt als Werkstudent in dem Unternehmen tätig. Da er dort auch sein eigenes Praxisprojekt absolviert hatte, 
konnte er mir detaillierte Einblicke in die technologische Ausrichtung und die Arbeitsabläufe geben. Auf seine Empfehlung hin initiierte ich den Bewerbungsprozess. 
In den darauffolgenden Gesprächen wurde sichergestellt, dass die Aufgabenstellung den Anforderungen des Fachbereichs Wirtschaftsinformatik entspricht und eine fachkundige Betreuung 
vor Ort gewährleistet ist. Da das Unternehmen moderne Frameworks und agile Methoden einsetzt, bot es das ideale Umfeld, um meine im Studium erworbenen theoretischen Kenntnisse in 
der Softwareentwicklung erstmals in einem professionellen, rein entwicklungsorientierten Kontext anzuwenden.